
\subsection*{7. Appendice~: Th\'eor\`eme de l'indice de Hodge}

On retrouve et compl\`ete ici quelques r\'esultats ant\'erieurs par un raisonnement direct.

\textbf{Th\'eor\`eme 7.1.} \emph{Soit @@M0@@ une surface irr\'eductible, propre sur un corps alg\'ebriquement clos. On suppose qu'il existe un faisceau inversible @@M1@@ sur @@M2@@ tel que @@M3@@ (par exemple, @@M4@@ ample). Alors un faisceau inversible @@M5@@ sur @@M6@@ tel que}
@@M7@@
\emph{est num\'eriquement \'equivalent \`a @@M8@@. (Par cons\'equent, la forme d'intersection sur @@M9@@ (@@M10@@, d'apr\`es 4.6) est non-d\'eg\'en\'er\'ee de type @@M11@@).}

\textbf{7.1.1.} D'abord supposons : @@M12@@ r\'eduit (donc, int\`egre) ; @@M13@@ ample ; @@M14@@ ; et @@M15@@. Raisonnant par l'absurde, prenons @@M16@@ tel que @@M17@@ soit ample, et prenons @@M18@@ tel que, en posant @@M19@@, on ait que @@M20@@ est @@M21@@. Alors on a aussi que @@M22@@ est @@M23@@, et que @@M24@@ est @@M25@@, ce qui contredit le lemme suivant.

\textbf{Lemme 7.1.2.} \emph{Soit @@M26@@ une surface int\`egre, projective sur un corps alg\'ebriquement clos. Pour un faisceau inversible @@M27@@ sur @@M28@@ tel que @@M29@@, les conditions suivantes sont \'equivalentes}~:
\begin{enumerate}[label=(\roman*)]
\item \emph{Pour tout faisceau ample @@M30@@, on a @@M31@@.}
\item[(i')] \emph{Il existe un faisceau ample @@M32@@ tel qu'on ait @@M33@@.}
\item \emph{Il existe un entier @@M34@@ tel qu'on ait @@M35@@.}
\end{enumerate}

